\documentclass[11pt,a4paper]{article}
\usepackage[utf8]{inputenc}
\usepackage[T1]{fontenc}
\usepackage[spanish,es-minimal,es-noshorthands]{babel}
\usepackage{lmodern}
\linespread{1.22}
\usepackage[a4paper,top=2.4cm,bottom=2.2cm,left=1.5cm,right=1.5cm,headheight=1.1cm,headsep=0.5cm,footskip=1cm]{geometry}
\usepackage{amsmath,amssymb,cancel}
\usepackage{icomma}
\usepackage{graphicx}
\usepackage[table]{xcolor}
\usepackage{tikz}\usetikzlibrary{calc}
\usepackage{pgfplots}\pgfplotsset{compat=1.17}
\usepackage{fontawesome5}
\usepackage{ccicons}
\usepackage{fancyhdr}
\usepackage{enumitem}
\usepackage{booktabs}
\usepackage{tabularx}
\usepackage{needspace}
\usepackage[most]{tcolorbox}
\usepackage{titlesec}
\usepackage[hidelinks]{hyperref}
\hypersetup{pdftitle={Apuntes de teoría. Unidad 1: Números reales. 1.º de Bachillerato CCSS},pdfsubject={Matemáticas Aplicadas a las Ciencias Sociales I},pdfkeywords={números reales, fracción generatriz, valor absoluto, intervalos, potencias, radicales, racionalización, logaritmos, notación científica}}

% colors de la web (FentMates, bloque de Aritmética: amarillo; es manté el nom fmgreen dels apunts de les altres unitats)
\definecolor{fmnavy}{HTML}{1B2748}
\definecolor{fmgreen}{HTML}{C99A2E}
\definecolor{fmgreenl}{HTML}{E6B84A}
\definecolor{fmyellow}{HTML}{FFC94D}
\definecolor{fmgold}{HTML}{E9A919}
\definecolor{fmaccent}{HTML}{FF7A50}
\definecolor{fmteal}{HTML}{2BB3A3}
\definecolor{fmdef}{HTML}{1D4F73}\definecolor{fmdefbg}{HTML}{EEF4F9}
\definecolor{fmprop}{HTML}{3B4252}\definecolor{fmpropbg}{HTML}{F3F4F7}
\definecolor{fmex}{HTML}{A3203A}\definecolor{fmexbg}{HTML}{FBF0F2}
\definecolor{fmmuted}{HTML}{5D6685}\definecolor{fmball0}{HTML}{E2483D}\definecolor{fmball1}{HTML}{2F6FD6}
\definecolor{fmsurface}{HTML}{EEF1F8}
\definecolor{fmA}{HTML}{3B82F6}\definecolor{fmB}{HTML}{F97316}\definecolor{fmR}{HTML}{E11D48}\definecolor{fmK}{HTML}{16A34A}\definecolor{fmV}{HTML}{7C3AED}

% blocs que no es poden partir: es mesuren i es reserva exactament el seu espai (amb el títol o la capçalera del requadre)
\newsavebox{\fmkb}
\newcommand{\fmmeasure}[2]{\setbox\fmkb=\vbox{\hsize=\dimexpr\linewidth-#1\relax\linewidth=\hsize #2}}
\newcommand{\fmblock}[1]{\fmmeasure{0pt}{#1}\Needspace{\dimexpr\ht\fmkb+\dp\fmkb\relax}#1}
\newcommand{\fmheadblock}[2]{\fmmeasure{0pt}{#2}\Needspace{\dimexpr\ht\fmkb+\dp\fmkb+2.2cm\relax}#1#2}
\newcommand{\fmboxneed}[1]{\fmmeasure{9mm}{#1}\Needspace{\dimexpr\ht\fmkb+\dp\fmkb+1.6cm\relax}}
\newif\iffmstack\newif\iffmfit
\newcommand{\fmfits}[1]{\par\ifdim\dimexpr\pagegoal-\pagetotal\relax<\dimexpr\ht\fmkb+\dp\fmkb+#1\relax\global\fmfitfalse\else\global\fmfittrue\fi}
\newcommand{\fmtwo}[2]{\fmmeasure{0pt}{#1}\fmfits{0pt}\iffmfit #1\else #2\fi}
\newcommand{\fmheadtwo}[3]{\fmmeasure{0pt}{#2}\fmfits{2.2cm}\iffmfit #1#2\else\Needspace{\dimexpr1.6cm+7\baselineskip\relax}#1#3\fi}
\newcommand{\fmtwoin}[2]{\iffmstack #2\else #1\fi}
\newcommand{\fmboxneedalt}[1]{\fmmeasure{9mm}{#1}\fmfits{1.6cm}\iffmfit\global\fmstackfalse\else\global\fmstacktrue\Needspace{8\baselineskip}\fi}
\newcommand{\fmheadboxalt}[2]{\fmmeasure{9mm}{#2}\fmfits{3.4cm}\iffmfit\global\fmstackfalse\else\global\fmstacktrue\Needspace{\dimexpr2.4cm+5\baselineskip\relax}\fi#1}
\newcommand{\fmheadbox}[2]{\fmmeasure{9mm}{#2}\Needspace{\dimexpr\ht\fmkb+\dp\fmkb+3.4cm\relax}#1}
\newcommand{\lt}{<}\newcommand{\gt}{>}
% fórmula en línia que no es parteix mai; si no cap en una línia, es posa centrada
% fórmula en línia; si ocupa més del 80 % de la línia, va centrada (amb la puntuació que la segueix)
\makeatletter
\newcommand{\fmm}[1]{\setbox0=\hbox{$#1$}\ifdim\wd0>0.8\linewidth\def\fmmnext{\fmmdisp{#1}}\else\def\fmmnext{\mbox{\box0}}\fi\fmmnext}
\def\fmmdisp#1{\@ifnextchar.{\fmmdispp{#1}}{\@ifnextchar,{\fmmdispc{#1}}{\@ifnextchar;{\fmmdisps{#1}}{\fmdisp{#1}}}}}
\def\fmmdispp#1.{\fmdisp{#1\,.}}\def\fmmdispc#1,{\fmdisp{#1\,,}}\def\fmmdisps#1;{\fmdisp{#1\,;}}
\makeatother
% fórmula centrada; si és més ampla que la línia, s'encongeix fins que hi cap
\newcommand{\fmdisp}[1]{\[\setbox2=\hbox{$\displaystyle #1$}\ifdim\wd2>\linewidth\resizebox{\linewidth}{!}{\box2}\else\box2\fi\]}
\DeclareUnicodeCharacter{2212}{-}
\DeclareUnicodeCharacter{2225}{\ensuremath{\parallel}}\DeclareUnicodeCharacter{22A5}{\ensuremath{\perp}}\DeclareUnicodeCharacter{2226}{\ensuremath{\nparallel}}
\DeclareUnicodeCharacter{207A}{\ensuremath{^{+}}}\DeclareUnicodeCharacter{207B}{\ensuremath{^{-}}}
\setlength{\parindent}{0pt}
\setlength{\parskip}{12pt}
\setlength{\lineskip}{5pt}\setlength{\lineskiplimit}{5pt}  % més aire al voltant de les fraccions en línia
\setlength{\emergencystretch}{3em}
\setlist{itemsep=4pt,topsep=4pt}
\renewcommand{\thesection}{1.\arabic{section}}
\titleformat{\section}{\color{fmnavy}\Large\bfseries}{\colorbox{fmdef}{\color{white}\strut\ \thesection\ }}{0.6em}{}[{\color{fmdef}\titlerule[1.2pt]}]
\titlespacing*{\section}{0pt}{10pt}{4pt}
\titleformat*{\subsection}{\color{fmgreen}\large\bfseries}
\titlespacing*{\subsection}{0pt}{6pt plus 2pt}{0pt}
\allowdisplaybreaks

% dibuix de la targeta de la unitat a la web (illustrations.js), en blanc
\newcommand{\fmcard}{\begin{tikzpicture}[x=0.062cm,y=-0.062cm,white,line width=3.2pt,line cap=round]
  \begin{scope}[transparency group,opacity=0.15]
  \draw (14,80)--(226,80) (216,72)--(228,80)--(216,88);
  \draw[line width=2.4pt] (40,74)--(40,86) (100,74)--(100,86) (160,74)--(160,86);
  \foreach \x/\t in {40/0,100/1,160/2} \node[anchor=base,inner sep=0pt,font=\fontsize{18}{18}\sffamily\bfseries] at (\x,106) {\t};
  \draw[line width=2.2pt,opacity=0.6] (40,80)--(40,20)--(100,20);
  \draw[line width=2.4pt] (40,80)--(100,20);
  \draw[dash pattern=on 6pt off 6pt] (100,20) arc[start angle=45,end angle=0,radius=84.85];
  \fill (125,80) circle (6.5) (194,80) circle (5);
  \node[anchor=base west,inner sep=0pt,font=\fontsize{20}{20}\selectfont] at (108,62) {$\sqrt2$};
  \node[anchor=base west,inner sep=0pt,font=\fontsize{20}{20}\selectfont] at (186,50) {$\pi$};
  \end{scope}
\end{tikzpicture}}
\newcommand{\fmbrand}{\raisebox{-0.25\height}{\includegraphics[height=0.5cm]{logo-fentmates.pdf}}}

% capçalera i peu
\pagestyle{fancy}\fancyhf{}
\fancyhead[L]{\small\fmbrand}
\fancyhead[R]{\small\color{fmmuted}1.º Bachillerato CCSS · Unidad 1: Números reales}
\fancyfoot[L]{\scriptsize\color{fmmuted}\href{https://creativecommons.org/licenses/by-nc-sa/4.0/deed.ca}{\ccbyncsa}\ \ FentMates · fentmates.org · CC BY-NC-SA 4.0}
\fancyfoot[R]{\small\color{fmmuted}\thepage}
\renewcommand{\headrulewidth}{0.4pt}\renewcommand{\headrule}{\color{fmmuted!40}\hrule height \headrulewidth}
\fancypagestyle{primera}{\fancyhf{}\renewcommand{\headrulewidth}{0pt}
  \fancyfoot[L]{\scriptsize\color{fmmuted}\href{https://creativecommons.org/licenses/by-nc-sa/4.0/deed.ca}{\ccbyncsa}\ \ FentMates · fentmates.org · CC BY-NC-SA 4.0}
  \fancyfoot[R]{\small\color{fmmuted}\thepage}}

% caixes (mateixos colors que la web)
% quan una caixa es parteix entre dues pàgines, la vora del tall és discontínua
\tcbset{fmbase/.style={enhanced,breakable,arc=3mm,boxrule=1pt,left=3.5mm,right=3.5mm,top=2mm,bottom=2.5mm,before skip=10pt,after skip=10pt,fonttitle=\bfseries,
  frame style={fill=tcbcolframe},
  overlay first={\fill[tcbcolback] ([xshift=1pt]frame.south west) rectangle ([xshift=-1pt,yshift=1.2pt]frame.south east);
    \draw[tcbcolframe,line width=1pt,dash pattern=on 5pt off 3pt] (frame.south west)--(frame.south east);},
  overlay middle={\fill[tcbcolback] ([xshift=1pt]frame.south west) rectangle ([xshift=-1pt,yshift=1.2pt]frame.south east);
    \fill[tcbcolback] ([xshift=1pt]frame.north west) rectangle ([xshift=-1pt,yshift=-1.2pt]frame.north east);
    \draw[tcbcolframe,line width=1pt,dash pattern=on 5pt off 3pt] (frame.south west)--(frame.south east) (frame.north west)--(frame.north east);},
  overlay last={\fill[tcbcolback] ([xshift=1pt]frame.north west) rectangle ([xshift=-1pt,yshift=-1.2pt]frame.north east);
    \draw[tcbcolframe,line width=1pt,dash pattern=on 5pt off 3pt] (frame.north west)--(frame.north east);}}}
\newtcolorbox{fmdef}[1]{fmbase,colback=fmdefbg,colframe=fmdef,coltitle=white,title={#1}}
\newtcolorbox{fmprop}[1]{fmbase,colback=fmpropbg,colframe=fmprop,coltitle=white,title={#1},halign title=center}
\newtcolorbox{fmex}[1]{fmbase,before upper=\raggedright,colback=fmexbg,colframe=fmex,colbacktitle=fmexbg,coltitle=fmex,title={\faPencil*\ #1},
  titlerule=0.4pt,titlerule style={fmex!40,dashed}}
\newtcolorbox{fmwarn}[1]{fmbase,colback=fmsf!6,colframe=fmsf,coltitle=white,title={#1}}
\newtcolorbox{fmnote}[1][]{enhanced,breakable,colback=white,colframe=white,borderline west={3pt}{0pt}{fmaccent},arc=0pt,left=3mm,right=2mm,top=1mm,bottom=1mm,before skip=6pt,after skip=6pt,before upper=\raggedright,#1}
\newcommand{\fmkey}[1]{\colorbox{fmyellow!45}{\strut #1}}
\newcommand{\fmuline}[1]{\textbf{\color{fmaccent}#1}}
\definecolor{fmok}{HTML}{15803D}\definecolor{fmev}{HTML}{2563EB}\definecolor{fmsf}{HTML}{C2410C}\definecolor{fmsi}{HTML}{7C3AED}\definecolor{fmas}{HTML}{A3203A}
\newcommand{\fmbadge}[2]{\tcbox[on line,colback=#1,colframe=#1,arc=2mm,boxsep=0pt,left=4pt,right=4pt,top=1.5pt,bottom=1.5pt]{\color{white}\footnotesize\bfseries #2}}
\newcommand{\fmchip}[1]{\tcbox[on line,colback=fmsf!10,colframe=fmsf!50,arc=2mm,boxsep=1pt,left=4pt,right=4pt]{#1}}
\newcommand{\fmlevel}[2]{\tcbox[on line,colback=white,colframe=fmmuted!40,arc=2mm,boxsep=2pt,halign=center]{\begin{tabular}{c}{\scriptsize\color{fmmuted}\MakeUppercase{#1}}\\#2\end{tabular}}}
% gràfiques
\pgfplotsset{fmaxis/.style={axis lines=middle,grid=major,grid style={fmmuted!15},axis line style={fmmuted,-latex},
  tick label style={font=\scriptsize,color=fmmuted},clip=true,unbounded coords=jump,restrict y to domain=-200:200,
  every axis plot/.append style={line join=round}},
  fmcurve/.style={fmaccent,line width=1.6pt,smooth}}
\tikzset{fmasy/.style={fmmuted,dashed,line width=0.9pt},fmguide/.style={fmaccent,dotted,line width=0.8pt},
  fmdot/.style={circle,fill=fmaccent,inner sep=0pt,minimum size=5pt},fmopen/.style={circle,draw=fmaccent,fill=white,line width=1pt,inner sep=0pt,minimum size=5pt}}
\newcommand{\fmpau}[1]{\tcbox[on line,colback=fmnavy,colframe=fmnavy,arc=2mm,boxsep=0pt,left=4pt,right=4pt,top=1.5pt,bottom=1.5pt]{\color{white}\scriptsize\bfseries\MakeUppercase{#1}}\ }

\begin{document}
\thispagestyle{primera}
\begin{tikzpicture}[remember picture,overlay]
  \fill[fmgreen] (current page.north west) rectangle ([yshift=-10.1cm]current page.north east);
  \begin{scope}
    \clip (current page.north west) rectangle ([yshift=-10.1cm]current page.north east);
    \node[anchor=north west,inner sep=0pt] at ([xshift=6.6cm,yshift=-1.3cm]current page.north west) {\fmcard};
  \end{scope}
  \node[anchor=north west,inner sep=0pt] at ([xshift=1.5cm,yshift=-1.1cm]current page.north west) {\tcbox[colback=white,colframe=white,arc=2.5mm,boxsep=0pt,left=7pt,right=7pt,top=5pt,bottom=5pt]{\includegraphics[height=0.75cm]{logo-fentmates.pdf}}};
\end{tikzpicture}
\vspace*{0.55cm}
\tcbox[colback=white!15!fmgreen,colframe=white!15!fmgreen,arc=2.5mm,boxsep=0pt,left=6pt,right=6pt,top=3pt,bottom=3pt]{\color{white}\small\bfseries Apuntes de teoría · 1.º Bachillerato · Matemáticas Aplicadas a las Ciencias Sociales I}

\vspace{4pt}
{\raggedright\color{white}\fontsize{25}{30}\selectfont\bfseries{\color{fmnavy}Unidad 1} · Números reales\par}
\vspace{6pt}
\begin{minipage}{\linewidth}\color{white!88!fmgreen}\small Teoría y ejemplos resueltos paso a paso: los conjuntos numéricos \fmm{\mathbb N,\ \mathbb Z,\ \mathbb Q,\ \mathbb I} y \fmm{\mathbb R}, los intervalos y las semirrectas, el valor absoluto, la unión y la intersección de intervalos, los radicales (forma exponencial, simplificación, índice común, extracción e introducción de factores y operaciones), la racionalización y los logaritmos con sus propiedades.\end{minipage}\par
% l'índex comença sempre a la mateixa altura, per davall de la franja (com en les unitats 5-10)
\vspace*{\dimexpr 8.6cm-\pagetotal\relax}

\begin{tcolorbox}[colback=fmsurface,colframe=fmsurface,arc=3mm,left=4mm,right=4mm,top=2mm,bottom=2mm]
\small\textbf{Índice}\par\smallskip
\begin{minipage}[t]{0.49\linewidth}\small \mbox{\hyperlink{sec1}{\textbf{\color{fmgreen}1.1}\ \ El conjunto de los números reales}}\\[3pt]
\mbox{\hyperlink{sec2}{\textbf{\color{fmgreen}1.2}\ \ Intervalos y semirrectas}}\\[3pt]
\mbox{\hyperlink{sec3}{\textbf{\color{fmgreen}1.3}\ \ Valor absoluto}}\\[3pt]
\mbox{\hyperlink{sec4}{\textbf{\color{fmgreen}1.4}\ \ Unión e intersección de intervalos}}\\[3pt]\end{minipage}\hfill\begin{minipage}[t]{0.49\linewidth}\small \mbox{\hyperlink{sec5}{\textbf{\color{fmgreen}1.5}\ \ Radicales}}\\[3pt]
\mbox{\hyperlink{sec6}{\textbf{\color{fmgreen}1.6}\ \ Racionalización}}\\[3pt]
\mbox{\hyperlink{sec7}{\textbf{\color{fmgreen}1.7}\ \ Logaritmos}}\\[3pt]
\mbox{\hyperlink{sec8}{\textbf{\color{fmgreen}1.8}\ \ Propiedades de los logaritmos}}\\[3pt]\end{minipage}
\end{tcolorbox}

\Needspace{8\baselineskip}
\section{El conjunto de los números reales}\hypertarget{sec1}{}

Recordemos que el conjunto de los \textbf{números racionales} \fmm{(\mathbb Q)} está formado por todos aquellos números que podemos escribir como una fracción \fmm{\dfrac ab}, con \fmm{a,b\in\mathbb Z} y \fmm{b\neq0}. El conjunto de los \textbf{números irracionales} \fmm{(\mathbb I)} está formado por los decimales con infinitas cifras decimales y no periódicos, por ejemplo \fmm{\pi,\ e,\ \sqrt2,\ \sqrt3} o el número de oro \fmm{\phi=\frac{1+\sqrt5}{2}}.

\Needspace{5\baselineskip}
\begin{fmdef}{Definición: números reales}
\begin{itemize}
\item \textbf{Naturales} \fmm{\mathbb N=\{0,1,2,3,\dots\}}: sirven para contar.
\item \textbf{Enteros} \fmm{\mathbb Z=\{\dots,-2,-1,0,1,2,\dots\}}: los naturales y sus opuestos.
\item \textbf{Racionales} \fmm{\mathbb Q}: los que se pueden escribir como fracción; su expresión decimal es exacta o periódica.
\item \textbf{Irracionales} \fmm{\mathbb I}: infinitas cifras decimales no periódicas (\fmm{\pi,\ \sqrt2,\ e\dots}).
\end{itemize}

El conjunto de los \textbf{números reales} \fmm{\mathbb R} es la unión del conjunto de los números racionales y el de los irracionales:

\fmdisp{\mathbb R=\mathbb Q\cup\mathbb I\qquad\qquad \mathbb N\subset\mathbb Z\subset\mathbb Q\subset\mathbb R,\qquad \mathbb I\subset\mathbb R,\qquad \mathbb Q\cap\mathbb I=\varnothing}

\end{fmdef}

\Needspace{5.1cm}
\subsection*{Expresiones decimales y fracción generatriz}
\begin{fmdef}{Tipos de decimales}
Un número es racional si y solo si su expresión decimal es \textbf{exacta} o \textbf{periódica}. Todos estos decimales tienen una \textbf{fracción generatriz}:

\begin{center}\renewcommand{\arraystretch}{1.5}\setlength{\parskip}{0pt}\small\begin{tabularx}{\linewidth}{>{\raggedright\arraybackslash}p{3.2cm}>{\raggedright\arraybackslash}X}\toprule
\textbf{Tipo} & \textbf{Ejemplo}\\\midrule
Exacto & $2{,}45=\dfrac{245}{100}=\dfrac{49}{20}$\\[3pt]\end{tabularx}\par\nointerlineskip\vskip0pt
\small\begin{tabularx}{\linewidth}{>{\raggedright\arraybackslash}p{3.2cm}>{\raggedright\arraybackslash}X}Periódico puro & $2{,}\widehat7=\dfrac{27-2}{9}=\dfrac{25}{9}$\\[3pt]\end{tabularx}\par\nointerlineskip\vskip0pt
\small\begin{tabularx}{\linewidth}{>{\raggedright\arraybackslash}p{3.2cm}>{\raggedright\arraybackslash}X}Periódico mixto & $1{,}2\widehat3=\dfrac{123-12}{90}=\dfrac{111}{90}=\dfrac{37}{30}$\\[3pt]\bottomrule\end{tabularx}\end{center}

\begin{fmnote}\fmuline{Regla}: en el numerador, el número sin coma (hasta el final del primer periodo) menos la parte que no se repite; en el denominador, tantos \textbf{9} como cifras tiene el periodo seguidos de tantos \textbf{0} como cifras tiene el anteperiodo.\end{fmnote}
\end{fmdef}

\Needspace{8\baselineskip}
\section{Intervalos y semirrectas}\hypertarget{sec2}{}

Los números reales se pueden representar en una recta en la que cada punto es un número real: el conjunto de los reales es \textbf{completo}, no hay «huecos» en la recta. Como entre dos números siempre hay infinitos, para describir una parte de la recta usamos los \textbf{intervalos}. Si un extremo pertenece al conjunto, el intervalo es \textbf{cerrado} en ese lado y escribimos \fmm{[\ ]}; si no pertenece, es \textbf{abierto} y escribimos \fmm{(\ )}.

\Needspace{5\baselineskip}
\begin{fmdef}{Intervalos y semirrectas}
\begin{center}\renewcommand{\arraystretch}{1.5}\setlength{\parskip}{0pt}\small\begin{tabularx}{\linewidth}{>{\raggedright\arraybackslash}p{3.2cm}>{\raggedright\arraybackslash}X>{\raggedright\arraybackslash}X>{\raggedright\arraybackslash}X}\toprule
\textbf{Nombre} & \textbf{Intervalo} & \textbf{Desigualdad} & \textbf{Representación}\\\midrule
Cerrado & $[a,b]$ & $a\le x\le b$ & \begin{tikzpicture}[x=0.0125cm,y=-0.0125cm,baseline=-0.2cm]\draw[fmmuted,line width=0.8pt] (10,14)--(186,14);\fill[fmmuted] (192,14)--++(-8,-4)--++(0,8)--cycle;\draw[fmR,line width=2.2pt] (64,14)--(136,14);\draw[fmR,line width=1.1pt,fill=fmR] (64,14) circle[radius=4.5];\node[font=\footnotesize] at (64,29) {$a$};\draw[fmR,line width=1.1pt,fill=fmR] (136,14) circle[radius=4.5];\node[font=\footnotesize] at (136,29) {$b$};\end{tikzpicture}\\[3pt]\end{tabularx}\par\nointerlineskip\vskip0pt
\small\begin{tabularx}{\linewidth}{>{\raggedright\arraybackslash}p{3.2cm}>{\raggedright\arraybackslash}X>{\raggedright\arraybackslash}X>{\raggedright\arraybackslash}X}Abierto & $(a,b)$ & $a\lt x\lt b$ & \begin{tikzpicture}[x=0.0125cm,y=-0.0125cm,baseline=-0.2cm]\draw[fmmuted,line width=0.8pt] (10,14)--(186,14);\fill[fmmuted] (192,14)--++(-8,-4)--++(0,8)--cycle;\draw[fmR,line width=2.2pt] (64,14)--(136,14);\draw[fmR,line width=1.1pt,fill=white] (64,14) circle[radius=4.5];\node[font=\footnotesize] at (64,29) {$a$};\draw[fmR,line width=1.1pt,fill=white] (136,14) circle[radius=4.5];\node[font=\footnotesize] at (136,29) {$b$};\end{tikzpicture}\\[3pt]\end{tabularx}\par\nointerlineskip\vskip0pt
\small\begin{tabularx}{\linewidth}{>{\raggedright\arraybackslash}p{3.2cm}>{\raggedright\arraybackslash}X>{\raggedright\arraybackslash}X>{\raggedright\arraybackslash}X}Abierto por la derecha & $[a,b)$ & $a\le x\lt b$ & \begin{tikzpicture}[x=0.0125cm,y=-0.0125cm,baseline=-0.2cm]\draw[fmmuted,line width=0.8pt] (10,14)--(186,14);\fill[fmmuted] (192,14)--++(-8,-4)--++(0,8)--cycle;\draw[fmR,line width=2.2pt] (64,14)--(136,14);\draw[fmR,line width=1.1pt,fill=fmR] (64,14) circle[radius=4.5];\node[font=\footnotesize] at (64,29) {$a$};\draw[fmR,line width=1.1pt,fill=white] (136,14) circle[radius=4.5];\node[font=\footnotesize] at (136,29) {$b$};\end{tikzpicture}\\[3pt]\end{tabularx}\par\nointerlineskip\vskip0pt
\small\begin{tabularx}{\linewidth}{>{\raggedright\arraybackslash}p{3.2cm}>{\raggedright\arraybackslash}X>{\raggedright\arraybackslash}X>{\raggedright\arraybackslash}X}Abierto por la izquierda & $(a,b]$ & $a\lt x\le b$ & \begin{tikzpicture}[x=0.0125cm,y=-0.0125cm,baseline=-0.2cm]\draw[fmmuted,line width=0.8pt] (10,14)--(186,14);\fill[fmmuted] (192,14)--++(-8,-4)--++(0,8)--cycle;\draw[fmR,line width=2.2pt] (64,14)--(136,14);\draw[fmR,line width=1.1pt,fill=white] (64,14) circle[radius=4.5];\node[font=\footnotesize] at (64,29) {$a$};\draw[fmR,line width=1.1pt,fill=fmR] (136,14) circle[radius=4.5];\node[font=\footnotesize] at (136,29) {$b$};\end{tikzpicture}\\[3pt]\end{tabularx}\par\nointerlineskip\vskip0pt
\small\begin{tabularx}{\linewidth}{>{\raggedright\arraybackslash}p{3.2cm}>{\raggedright\arraybackslash}X>{\raggedright\arraybackslash}X>{\raggedright\arraybackslash}X}Semirrecta cerrada (derecha) & $[a,+\infty)$ & $x\ge a$ & \begin{tikzpicture}[x=0.0125cm,y=-0.0125cm,baseline=-0.2cm]\draw[fmmuted,line width=0.8pt] (10,14)--(186,14);\fill[fmmuted] (192,14)--++(-8,-4)--++(0,8)--cycle;\draw[fmR,line width=2.2pt] (100,14)--(182,14);\fill[fmR] (188,14)--++(-9,-6)--++(0,12)--cycle;\draw[fmR,line width=1.1pt,fill=fmR] (100,14) circle[radius=4.5];\node[font=\footnotesize] at (100,29) {$a$};\end{tikzpicture}\\[3pt]\end{tabularx}\par\nointerlineskip\vskip0pt
\small\begin{tabularx}{\linewidth}{>{\raggedright\arraybackslash}p{3.2cm}>{\raggedright\arraybackslash}X>{\raggedright\arraybackslash}X>{\raggedright\arraybackslash}X}Semirrecta abierta (derecha) & $(a,+\infty)$ & $x\gt a$ & \begin{tikzpicture}[x=0.0125cm,y=-0.0125cm,baseline=-0.2cm]\draw[fmmuted,line width=0.8pt] (10,14)--(186,14);\fill[fmmuted] (192,14)--++(-8,-4)--++(0,8)--cycle;\draw[fmR,line width=2.2pt] (100,14)--(182,14);\fill[fmR] (188,14)--++(-9,-6)--++(0,12)--cycle;\draw[fmR,line width=1.1pt,fill=white] (100,14) circle[radius=4.5];\node[font=\footnotesize] at (100,29) {$a$};\end{tikzpicture}\\[3pt]\end{tabularx}\par\nointerlineskip\vskip0pt
\small\begin{tabularx}{\linewidth}{>{\raggedright\arraybackslash}p{3.2cm}>{\raggedright\arraybackslash}X>{\raggedright\arraybackslash}X>{\raggedright\arraybackslash}X}Semirrecta cerrada (izquierda) & $(-\infty,b]$ & $x\le b$ & \begin{tikzpicture}[x=0.0125cm,y=-0.0125cm,baseline=-0.2cm]\draw[fmmuted,line width=0.8pt] (10,14)--(186,14);\fill[fmmuted] (192,14)--++(-8,-4)--++(0,8)--cycle;\draw[fmR,line width=2.2pt] (14,14)--(100,14);\fill[fmR] (8,14)--++(9,-6)--++(0,12)--cycle;\draw[fmR,line width=1.1pt,fill=fmR] (100,14) circle[radius=4.5];\node[font=\footnotesize] at (100,29) {$b$};\end{tikzpicture}\\[3pt]\end{tabularx}\par\nointerlineskip\vskip0pt
\small\begin{tabularx}{\linewidth}{>{\raggedright\arraybackslash}p{3.2cm}>{\raggedright\arraybackslash}X>{\raggedright\arraybackslash}X>{\raggedright\arraybackslash}X}Semirrecta abierta (izquierda) & $(-\infty,b)$ & $x\lt b$ & \begin{tikzpicture}[x=0.0125cm,y=-0.0125cm,baseline=-0.2cm]\draw[fmmuted,line width=0.8pt] (10,14)--(186,14);\fill[fmmuted] (192,14)--++(-8,-4)--++(0,8)--cycle;\draw[fmR,line width=2.2pt] (14,14)--(100,14);\fill[fmR] (8,14)--++(9,-6)--++(0,12)--cycle;\draw[fmR,line width=1.1pt,fill=white] (100,14) circle[radius=4.5];\node[font=\footnotesize] at (100,29) {$b$};\end{tikzpicture}\\[3pt]\bottomrule\end{tabularx}\end{center}

\begin{fmnote}Las \textbf{semirrectas} tienen un extremo en el infinito, y en el infinito el intervalo siempre es abierto. La recta entera es \fmm{\mathbb R=(-\infty,+\infty)}. En algunos libros el extremo abierto se escribe con el corchete girado: \fmm{]-2,2[\ =(-2,2)}.\end{fmnote}
\end{fmdef}

\Needspace{5\baselineskip}
\begin{fmex}{Ejemplo resuelto}
\begin{enumerate}[label=\textbf{\arabic*.},leftmargin=*]
\item Escribe como intervalo y representa «los números mayores que \fmm{-1} y menores o iguales que 3».

Desigualdad: \fmm{\ -1\lt x\le 3}

El \fmm{-1} no está (abierto) y el 3 sí (cerrado): \fmm{\ x\in(-1,\,3]}

En la recta: punto vacío en el \fmm{-1}, punto lleno en el 3 y el trozo entre ambos pintado.
\item Escribe como semirrecta «los números menores que 2».

\fmm{x\lt 2\ \Rightarrow\ x\in(-\infty,\,2)}: abierta en el 2 porque el 2 no cumple la condición.
\end{enumerate}
\end{fmex}

\Needspace{5.7cm}
\section{Valor absoluto}\hypertarget{sec3}{}
\begin{fmdef}{Definición: valor absoluto}
Dado \fmm{x} un número real cualquiera, definimos el \textbf{valor absoluto} de \fmm{x} como:

\fmdisp{|x|=\begin{cases}\ \ x & \text{si } x\ge0\\ -x & \text{si } x\lt0\end{cases}}

El valor absoluto de un número es la \textbf{distancia} de ese número hasta el cero; por eso nunca es negativo. Por ejemplo, \fmm{|-5|=5=|5|}: los dos están a 5 unidades del cero. En general, la distancia entre dos números \fmm{a} y \fmm{b} es \fmm{d(a,b)=|a-b|}.
\end{fmdef}

\Needspace{5\baselineskip}
\begin{fmprop}{Propiedades del valor absoluto}

\fmdisp{|x|\ge0\qquad |-x|=|x|\qquad |x\cdot y|=|x|\cdot|y|\qquad |x+y|\le|x|+|y|}

\begin{fmnote}Para quitar el valor absoluto hay que saber el signo de lo que hay dentro: \fmm{|\pi-4|=4-\pi} porque \fmm{\pi-4\lt0}.\end{fmnote}
\end{fmprop}

\Needspace{5\baselineskip}
\begingroup\tcbset{before skip=5pt,after skip=5pt}\setlist{itemsep=2pt,topsep=0pt}
\begin{fmex}{Ejemplo resuelto: ecuaciones con valor absoluto}
Si \fmm{r\gt0}, \fmm{\ |A|=r\iff A=r\ \ \text{o}\ \ A=-r}: hay dos números a distancia \fmm{r} del cero.

\begin{enumerate}[label=\textbf{\arabic*.},leftmargin=*]
\item Resuelve \fmm{|2x-3|=7}: \fmm{2x-3=7\ \Rightarrow\ x=5}\ \ o\ \ \fmm{2x-3=-7\ \Rightarrow\ x=-2}. Soluciones: \fmm{x=5} y \fmm{x=-2}.
\end{enumerate}

\begin{fmnote}\fmm{|A|=-3} no tiene solución: un valor absoluto nunca es negativo.\end{fmnote}
\end{fmex}

\Needspace{5\baselineskip}
\begin{fmdef}{Valor absoluto e intervalos}
Como \fmm{|x-a|} es la distancia de \fmm{x} a \fmm{a}, las desigualdades con valor absoluto describen intervalos de la recta: \fmm{|x-a|\lt r} son los puntos a distancia menor que \fmm{r} del centro \fmm{a}, es decir, el intervalo \fmm{(a-r,\ a+r)}.

\fmdisp{|A|\lt r\iff -r\lt A\lt r\qquad\qquad |A|\gt r\iff A\lt -r\ \ \text{o}\ \ A\gt r}

\end{fmdef}

\Needspace{5\baselineskip}
\begin{fmex}{Ejemplo resuelto: desigualdades con valor absoluto}
\begin{enumerate}[label=\textbf{\arabic*.},leftmargin=*]
\item Resuelve \fmm{|3x-6|\le9}: \fmm{-9\le 3x-6\le 9}. Sumamos 6 a los tres miembros: \fmm{\ -3\le 3x\le 15}. Dividimos entre 3: \fmm{\ -1\le x\le 5\ \Rightarrow\ x\in[-1,\,5]}
\item Resuelve \fmm{|-2x+4|\gt6}: \fmm{-2x+4\gt 6\ \ \text{o}\ \ -2x+4\lt -6\ \ \Rightarrow\ \ -2x\gt 2\ \ \text{o}\ \ -2x\lt -10}. Al dividir entre un número negativo, la desigualdad \textbf{cambia de sentido}: \fmm{\ x\lt -1\ \ \text{o}\ \ x\gt 5\ \Rightarrow\ x\in(-\infty,-1)\cup(5,+\infty)}
\end{enumerate}

\begin{fmnote}\fmuline{Cuidado}: \fmm{|A|\lt -2} no tiene solución (un valor absoluto nunca es negativo) y \fmm{|A|\gt -1} se cumple para todo \fmm{x\in\mathbb R}.\end{fmnote}
\end{fmex}

\enlargethispage{0.4cm}\Needspace{5.7cm}
\section{Unión e intersección de intervalos}\hypertarget{sec4}{}
\begin{fmdef}{Operaciones con intervalos}
Dados dos intervalos \fmm{A} y \fmm{B}:

\begin{itemize}
\item La \textbf{unión} \fmm{A\cup B} está formada por los números que están en \fmm{A} \textbf{o} en \fmm{B} (en alguno de los dos).
\item La \textbf{intersección} \fmm{A\cap B} está formada por los números que están en \fmm{A} \textbf{y} en \fmm{B} (en los dos a la vez).
\end{itemize}

\begin{fmnote}Lo más cómodo es dibujar los dos intervalos en la misma recta, uno encima del otro. Si no tienen ningún número en común, la intersección es el \textbf{conjunto vacío} \fmm{\varnothing}; y a veces la unión no se puede escribir como un único intervalo: \fmm{[0,2)\cup(2,4]}.\end{fmnote}
\end{fmdef}

\begin{fmex}{Ejemplo resuelto}
\begin{enumerate}[label=\textbf{\arabic*.},leftmargin=*]
\item Sean \fmm{A=[-3,2]}, \fmm{B=(0,5]} y \fmm{C=[4,+\infty)}. Calcula \fmm{A\cap B}, \fmm{A\cup B}, \fmm{B\cap C} y \fmm{A\cap C}. Dibujamos los tres intervalos en la misma recta: así se ve enseguida qué tienen en común.

\fmm{A\cap B=(0,\,2]}: el 0 no es de \fmm{B} y el 2 sí que es de los dos. \fmm{A\cup B=[-3,\,5]}\qquad \fmm{B\cap C=[4,\,5]}\qquad \fmm{A\cap C=\varnothing}: \fmm{A} se acaba en el 2 y \fmm{C} empieza en el 4.
\end{enumerate}
\end{fmex}
\endgroup

\Needspace{5.7cm}
\section{Radicales}\hypertarget{sec5}{}
\begin{fmdef}{Definición: raíz enésima}
Dados \fmm{a} un número real y \fmm{n} un número natural, la \textbf{raíz enésima} de \fmm{a} es:

\fmdisp{\sqrt[n]{a}=x\iff a=x^{n}}

El número \fmm{a} es el \textbf{radicando} y \fmm{n} el \textbf{índice} del radical. Si el índice es par y el radicando negativo, la raíz no existe en \fmm{\mathbb R}: \fmm{\sqrt{-4}\notin\mathbb R}, pero \fmm{\sqrt[3]{-8}=-2}.
\end{fmdef}

\fmblock{\par\noindent \begin{minipage}[t]{0.48\linewidth}\vspace{0pt}\raggedright \begin{fmprop}{Propiedades de las potencias}

\fmdisp{\begin{array}{l} a^n\cdot a^m=a^{n+m}\\[2pt] a^n : a^m=a^{n-m}\\[2pt] a^n\cdot b^n=(a\cdot b)^n\\[2pt] a^n : b^n=(a:b)^n\\[2pt] (a^n)^m=a^{n\cdot m}\\[2pt] a^{-n}=\dfrac{1}{a^n}\qquad a^0=1\end{array}}

\end{fmprop}\end{minipage}\hfill\begin{minipage}[t]{0.48\linewidth}\vspace{0pt}\raggedright \begin{fmprop}{Propiedades de los radicales}

\fmdisp{\begin{array}{l} \sqrt[n\cdot p]{a^{p}}=\sqrt[n]{a}\\[2pt] \left(\sqrt[n]{a}\right)^p=\sqrt[n]{a^p}\\[2pt] \sqrt[m]{\sqrt[n]{a}}=\sqrt[m\cdot n]{a}\\[2pt] \sqrt[n]{a}\cdot\sqrt[n]{b}=\sqrt[n]{a\cdot b}\\[2pt] \dfrac{\sqrt[n]{a}}{\sqrt[n]{b}}=\sqrt[n]{\dfrac ab}\end{array}}

\end{fmprop}\end{minipage}\par\medskip}

\Needspace{5\baselineskip}
\begin{fmprop}{Errores frecuentes}
\begin{itemize}
\item \fmm{-3^2=-9\ } pero \fmm{\ (-3)^2=9}: el exponente solo afecta a lo que tiene justo delante.
\item \fmm{\sqrt{9+16}=\sqrt{25}=5\neq3+4}: la raíz \textbf{no} se reparte en sumas ni restas.
\item \fmm{(a+b)^2=a^2+2ab+b^2\neq a^2+b^2}
\item \fmm{2^{-3}=\frac18\gt0}: un exponente negativo no hace negativo el resultado.
\end{itemize}
\end{fmprop}

\Needspace{5.1cm}
\subsection*{Forma exponencial y simplificación}
\begin{fmdef}{Radicales como potencias}
Todo radical se puede escribir como una potencia de exponente fraccionario: el exponente del radicando va al numerador y el índice, al denominador.

\fmdisp{\sqrt[n]{a^m}=a^{\frac mn}\qquad\quad \sqrt5=5^{\frac12},\quad \sqrt[3]{x^2}=x^{\frac23},\quad 7^{\frac34}=\sqrt[4]{7^3},\quad 2^{-\frac12}=\dfrac{1}{\sqrt2}}

Así, todas las propiedades de los radicales son las de las potencias. Para \textbf{simplificar} un radical, simplificamos la fracción del exponente (dividimos el índice y el exponente entre el mismo número):

\fmdisp{\sqrt[6]{2^4}=2^{\frac46}=2^{\frac23}=\sqrt[3]{2^2}=\sqrt[3]4\qquad\quad \sqrt[12]{a^8b^4}=\sqrt[3]{a^2b}}

\end{fmdef}

\Needspace{5.1cm}
\subsection*{Reducción a índice común}
\begin{fmdef}{Índice común}
Para multiplicar, dividir o comparar radicales de índice distinto, los escribimos con el mismo índice: el \textbf{m.c.m.} de los índices. Cada radicando se eleva al número por el que hemos multiplicado su índice.

\fmdisp{\sqrt{3},\ \sqrt[3]{5}\ \xrightarrow{\ \text{m.c.m.}(2,3)=6\ }\ \sqrt[6]{3^3}=\sqrt[6]{27},\quad \sqrt[6]{5^2}=\sqrt[6]{25}\ \Rightarrow\ \sqrt[3]5\lt\sqrt3}

\end{fmdef}

\Needspace{5.1cm}
\subsection*{Extracción e introducción de factores}
\begin{fmex}{Ejemplo resuelto}
\begin{enumerate}[label=\textbf{\arabic*.},leftmargin=*]
\item \textbf{Introducir} factores: entran en la raíz elevados al índice.

\fmm{2^4\cdot 5\sqrt[3]{2\cdot 7}=\sqrt[3]{(2^4)^3\cdot 5^3\cdot 2\cdot 7}}

\fmm{=\sqrt[3]{2^{13}\cdot 5^3\cdot 7}}
\item \textbf{Extraer} factores: solo salen los que tienen exponente mayor o igual que el índice. Dividimos cada exponente entre el índice: el cociente es el exponente de fuera y el resto el de dentro.

\fmm{\sqrt[4]{2^7\cdot 3^2\cdot 5^8}}: \fmm{\quad 7=1\cdot4+3,\quad 2\lt4,\quad 8=2\cdot4+0}

\fmm{=2\cdot 5^2\sqrt[4]{2^3\cdot 3^2}=50\sqrt[4]{72}}
\end{enumerate}
\end{fmex}

\Needspace{5.1cm}
\subsection*{Operaciones con radicales}
\begin{fmdef}{Suma, resta, producto, cociente y potencia}
\begin{itemize}
\item \textbf{Sumas y restas}: solo se pueden sumar radicales \textbf{semejantes} (mismo índice y mismo radicando): \fmm{2\sqrt[3]{21}+4\sqrt[3]{21}=6\sqrt[3]{21}}, pero \fmm{2\sqrt{21}+4\sqrt[3]{21}} no se puede simplificar. A veces hay que extraer factores para descubrirlo: \fmm{\sqrt8+\sqrt{18}=2\sqrt2+3\sqrt2=5\sqrt2}.
\item \textbf{Producto y cociente}: con el mismo índice, se multiplican o se dividen los radicandos: \fmm{\sqrt{12}\cdot\sqrt3=\sqrt{36}=6}, \fmm{\ \dfrac{\sqrt[3]{40}}{\sqrt[3]5}=\sqrt[3]8=2}. Si no tienen el mismo índice, primero los reducimos a índice común: \fmm{\sqrt[3]{5}\cdot\sqrt2=\sqrt[6]{5^2}\cdot\sqrt[6]{2^3}=\sqrt[6]{200}}.
\item \textbf{Potencia}: se eleva el radicando: \fmm{\left(\sqrt[3]{4}\right)^2=\sqrt[3]{16}=2\sqrt[3]2}, \fmm{\ \left(\sqrt3\right)^4=\sqrt{3^4}=9}.
\item \textbf{Raíz de una raíz}: se multiplican los índices: \fmm{\sqrt{\sqrt[3]{4}}=\sqrt[6]{4}=\sqrt[3]{2}}.
\end{itemize}
\end{fmdef}

\Needspace{5\baselineskip}
\begin{fmex}{Ejemplo resuelto}
\begin{enumerate}[label=\textbf{\arabic*.},leftmargin=*]
\item \fmm{3\sqrt{50}-2\sqrt{18}+\sqrt{8}}

Extraemos factores y sumamos: \fmm{\ 3\cdot5\sqrt2-2\cdot3\sqrt2+2\sqrt2=(15-6+2)\sqrt2=11\sqrt2}

\item \fmm{\dfrac{\sqrt[4]{8}\cdot\sqrt2}{\sqrt[3]{4}}}

Índice común: m.c.m.\fmm{(4,2,3)=12}:\\[2ex]

\fmm{\dfrac{\sqrt[12]{2^9}\cdot\sqrt[12]{2^6}}{\sqrt[12]{2^8}}=\sqrt[12]{2^{9+6-8}}=\sqrt[12]{2^7}=\sqrt[12]{128}}

Con potencias es más rápido: \fmm{\ 2^{\frac34+\frac12-\frac23}=2^{\frac{7}{12}}}
\end{enumerate}
\end{fmex}

\Needspace{5.7cm}
\section{Racionalización}\hypertarget{sec6}{}
\begin{fmdef}{Racionalizar}
\textbf{Racionalizar} es buscar una fracción equivalente sin raíces en el denominador. Hay dos casos:

\begin{itemize}
\item En el denominador hay una \textbf{única raíz}: multiplicamos arriba y abajo por lo que haga falta para que desaparezca.
\item En el denominador hay una \textbf{suma o resta con raíces}: multiplicamos por el \textbf{conjugado} y usamos \fmm{(a+b)(a-b)=a^2-b^2}.
\end{itemize}
\end{fmdef}

\Needspace{5\baselineskip}
\begin{fmex}{Ejemplo resuelto}
\begin{enumerate}[label=\textbf{\arabic*.},leftmargin=*]
\item \fmm{\dfrac{2}{\sqrt[3]{5}}=\dfrac{2\cdot\sqrt[3]{5^2}}{\sqrt[3]{5}\cdot\sqrt[3]{5^2}}=\dfrac{2\sqrt[3]{5^2}}{\sqrt[3]{5^3}}=\dfrac{2\sqrt[3]{25}}{5}}


\item \fmm{\dfrac{5}{1+\sqrt3}=\dfrac{5}{1+\sqrt3}\cdot\dfrac{1-\sqrt3}{1-\sqrt3}=\dfrac{5(1-\sqrt3)}{1^2-(\sqrt3)^2}=\dfrac{5(1-\sqrt3)}{-2}=-\dfrac52\left(1-\sqrt3\right)}


\item \fmm{\dfrac{2}{\sqrt2-\sqrt3}=\dfrac{2(\sqrt2+\sqrt3)}{(\sqrt2)^2-(\sqrt3)^2}=\dfrac{2(\sqrt2+\sqrt3)}{-1}=-2\left(\sqrt2+\sqrt3\right)}

\end{enumerate}
\end{fmex}

\Needspace{5.7cm}
\section{Logaritmos}\hypertarget{sec7}{}
\begin{fmdef}{Definición: logaritmo}
El \textbf{logaritmo en base \fmm{a} de \fmm{b}} es el exponente al que hay que elevar \fmm{a} para obtener \fmm{b}:

\fmdisp{\log_a b=c\iff a^c=b\qquad (a\gt0,\ a\neq1,\ b\gt0)}

Por ejemplo, \fmm{\log_{10}1000=3} porque \fmm{10^3=1000}. Cuando la base es 10 escribimos \fmm{\log} (logaritmo decimal) y cuando es el número \fmm{e} escribimos \fmm{\ln} (logaritmo neperiano). Como \fmm{a^0=1}, \fmm{\log_a1=0} en cualquier base.
\end{fmdef}

\Needspace{5\baselineskip}
\begin{fmprop}{Consecuencias de la definición}

\fmdisp{\log_a 1=0\qquad \log_a a=1\qquad \log_a a^n=n\qquad a^{\log_a b}=b}

\begin{fmnote}No existen los logaritmos de números negativos ni del cero: \fmm{\log(-10)} y \fmm{\log0} no son números reales, porque ninguna potencia de base positiva da un resultado negativo o cero.\end{fmnote}
\end{fmprop}

\Needspace{5\baselineskip}
\begin{fmex}{Ejemplo resuelto: logaritmos con la definición}
\begin{enumerate}[label=\textbf{\arabic*.},leftmargin=*]
\item \fmm{\log_2\dfrac18}. Buscamos el exponente: \fmm{\ 2^c=\dfrac18=2^{-3}\ \Rightarrow\ \log_2\dfrac18=-3}
\item \fmm{\log_3\sqrt[5]{9}}: \fmm{\sqrt[5]{9}=\sqrt[5]{3^2}=3^{\frac25}\ \Rightarrow\ \log_3 3^{\frac25}=\dfrac25}
\item Halla \fmm{x} si \fmm{\log_x 125=3}: \fmm{x^3=125=5^3\ \Rightarrow\ x=5}
\end{enumerate}
\end{fmex}

\fmheadbox{
\section{Propiedades de los logaritmos}\hypertarget{sec8}{}
}{\par\noindent \begin{minipage}[t]{0.48\linewidth}\vspace{0pt}\raggedright
\fmdisp{\begin{array}{l}\log_a a=1\\[2pt] \log_a a^n=n\\[2pt] \log_a(p\cdot q)=\log_a p+\log_a q\\[2pt] \log_a(p:q)=\log_a p-\log_a q\\[2pt] \log_a b^n=n\cdot\log_a b\\[2pt] \log_a x=\dfrac{\log_b x}{\log_b a}\end{array}}
\end{minipage}\hfill\begin{minipage}[t]{0.48\linewidth}\vspace{0pt}\raggedright \textbf{Ejemplos}:

\fmm{\log10+\log5=\log(10\cdot5)=\log50}

\fmm{\log10-\log5=\log\tfrac{10}{5}=\log2}

\fmm{\ln 5^3=3\ln5}

Cambio de base: \fmm{\ \log_3 20=\dfrac{\log20}{\log3}=\dfrac{\ln20}{\ln3}\approx 2{,}727}\end{minipage}\par\medskip}
\begin{fmprop}{Propiedades de los logaritmos}
\par\noindent \begin{minipage}[t]{0.48\linewidth}\vspace{0pt}\raggedright
\fmdisp{\begin{array}{l}\log_a a=1\\[2pt] \log_a a^n=n\\[2pt] \log_a(p\cdot q)=\log_a p+\log_a q\\[2pt] \log_a(p:q)=\log_a p-\log_a q\\[2pt] \log_a b^n=n\cdot\log_a b\\[2pt] \log_a x=\dfrac{\log_b x}{\log_b a}\end{array}}
\end{minipage}\hfill\begin{minipage}[t]{0.48\linewidth}\vspace{0pt}\raggedright \textbf{Ejemplos}:

\fmm{\log10+\log5=\log(10\cdot5)=\log50}

\fmm{\log10-\log5=\log\tfrac{10}{5}=\log2}

\fmm{\ln 5^3=3\ln5}

Cambio de base: \fmm{\ \log_3 20=\dfrac{\log20}{\log3}=\dfrac{\ln20}{\ln3}\approx 2{,}727}\end{minipage}\par\medskip

\begin{fmnote}\fmuline{Cuidado}: \fmm{\log(p+q)\neq\log p+\log q}. La propiedad es para el \textbf{producto}.\end{fmnote}
\end{fmprop}

\Needspace{5\baselineskip}
\begin{fmex}{Ejemplo resuelto: aplicar las propiedades}
\begin{enumerate}[label=\textbf{\arabic*.},leftmargin=*]
\item Calcula \fmm{\log_2 48-\log_2 3} y \fmm{\ 2\log 5+\log 4}.

\fmm{\log_2 48-\log_2 3=\log_2\dfrac{48}{3}=\log_2 16=4}

\fmm{2\log5+\log4=\log5^2+\log4=\log(25\cdot4)=\log100=2}
\item Si \fmm{\log2\approx 0{,}301} y \fmm{\log3\approx 0{,}477}, calcula \fmm{\log 12} y \fmm{\log 1{,}5}.

\fmm{\log12=\log(2^2\cdot3)=2\log2+\log3\approx 1{,}079}

\fmm{\log 1{,}5=\log\dfrac32=\log3-\log2\approx 0{,}176}
\item Desarrolla \fmm{\log\dfrac{100\,x^3}{\sqrt y}}.

\fmm{=\log100+\log x^3-\log y^{\frac12}}

\fmm{=2+3\log x-\dfrac12\log y}
\item Expresa \fmm{A} sin logaritmos si \fmm{\log A=2\log x-\log y+1}.

\fmm{\log A=\log x^2-\log y+\log10=\log\dfrac{10x^2}{y}}

\fmm{A=\dfrac{10x^2}{y}}
\item Un capital de 5000 € se coloca al 4 \% de interés compuesto anual. ¿Cuántos años hacen falta para que se convierta en 8000 €?

\fmm{5000\cdot 1{,}04^t=8000\ \Rightarrow\ 1{,}04^t=1{,}6}

Tomamos logaritmos: \fmm{\ t\log 1{,}04=\log 1{,}6\ \Rightarrow\ t=\dfrac{\log 1{,}6}{\log 1{,}04}\approx 11{,}98}

Hacen falta 12 años.
\end{enumerate}
\end{fmex}

  
\end{document}
